% !TeX root = ../main.tex
% ─────────────────────────────────────────────────────────────
%  PROVENANCE — Cờ Trận is a folk game; the author did NOT create it.
%  Everything here comes from the author's recollection (14 Oct 2026).
% ─────────────────────────────────────────────────────────────
\section{Origin and Recording Method}
\label{sec:provenance}

\subsection{Place, players and name}

Cờ Trận was a children's game in Đồng Bùa hamlet, Tam Quan commune, Tam Đảo
district, Vĩnh Phúc province, Vietnam. We use the place names of the
administrative boundaries at the time the game was played; since 2025,
Vĩnh Phúc province has been merged into Phú Thọ province. The author was
taught Cờ Trận in grade 5, around the 2011--2012 school year, and played it
with friends in the village until the end of grade 9, around 2016, and then
stopped. The game was thus passed on to the author by earlier players; it
did not start with the author's generation.

A four-player game usually lasted three to four hours, from midday until
the evening. The folk game has \emph{no} move limit: a game ends only by the
winning conditions in Section~\ref{sec:rules}.

``Cờ Trận'' is the name used in the village. \emph{Cờ} means a board game,
and \emph{trận} evokes a battle or a battle formation. The English name
\emph{Tactical Formation} is the author's rendering of this name for an
international readership; it is not a folk name.

\subsection{Board and pieces}

The game needs no manufactured equipment. The board was drawn directly on the
road. On a dirt road, players scratched it into the ground with a
stick; on a concrete road, they drew it with a stone, chalk or a piece of
brick so that the lines would last longer. Pieces were whatever could be
picked up by the roadside: stones, pebbles, dry or green twigs, leaves.
Each side used a different kind of object. To capture, a player removed the
captured side's object and put one of their own in exactly that spot. This
keeps the total number of pieces on the board unchanged, matching the
conservation invariant in Appendix~\ref{app:formal}. Using found objects also
explains why the game has left no physical artefacts behind.

\subsection{Recording method and reliability}

The rules in this paper were written down from the author's own
recollection. According to the author, the record is complete and faithful
to the folk game in all of the following: the two-layer vertex--centre
board, the jump move, the sandwich rule and the choice of a single capture
direction, the number of pieces and the starting layout, the winning
condition, and all three modes for 2, 3 and 4 players.

We have not yet been able to check this record against a second source.
Most people in the village have forgotten the game. Web searches in
Vietnamese and English (October 2026) found no document describing Cờ Trận
or any game with the same rule set. The record therefore rests on \emph{a
single informant}, a limitation we state explicitly in
Section~\ref{sec:limitations}. We invite anyone who knows Cờ Trận, or a
variant of it from another village, to contact the author so that the
record can be extended and cross-checked.

\subsection{What is folk tradition and what was added}

So that readers do not confuse the folk game with the contributions of this
paper, we separate the two clearly:
\begin{itemize}
  \item \textbf{Part of the folk game}: all the rules specified in
        Sections~\ref{sec:rules} and~\ref{sec:variants}.
  \item \textbf{Added by the author}: the English name; the formal
        mathematical model (Appendix~\ref{app:formal}); the software and
        interface (Section~\ref{sec:impl}); the fixed turn order in the
        software (Section~\ref{ssec:turn_order}); the entire AI
        (Section~\ref{sec:ai}); and the experiments
        (Section~\ref{sec:eval}).
\end{itemize}
